% ═══════════════════════════════════════════════════════════════
% SPANISCH ARBEITSBLATT - VERSION E: "Refined Hybrid"
% ═══════════════════════════════════════════════════════════════
% Thema: Objetos de clase & hay
% Klasse 7 | Unidad 2 | Mi instituto
% ═══════════════════════════════════════════════════════════════

\documentclass[10pt, a4paper]{article}

\usepackage[utf8]{inputenc}
\usepackage[T1]{fontenc}
\usepackage[ngerman]{babel}

% --- SCHRIFTEN (B-Prinzip) ---
\usepackage{palatino}
\usepackage{helvet}
\newcommand{\sanstext}[1]{{\fontfamily{phv}\selectfont #1}}

\usepackage{microtype}
\usepackage[margin=1.4cm, top=1.6cm, bottom=1.3cm]{geometry}
\usepackage{enumitem}
\usepackage{tabularx}
\usepackage{booktabs}
\usepackage{array}
\usepackage{multicol}
\usepackage{tikz}
\usetikzlibrary{calc}

\usepackage{fancyhdr}
\usepackage{lastpage}
\usepackage{needspace}

\usepackage{xcolor}
\usepackage{amssymb}
\usepackage{tcolorbox}
\tcbuselibrary{skins}

% ═══════════════════════════════════════════════════════════════
% FARBEN
% ═══════════════════════════════════════════════════════════════

\definecolor{kopfzeile}{HTML}{1A1A1A}
\definecolor{akzent}{HTML}{C41E3A}          % Spanisch-Karminrot
\definecolor{hellgrau}{HTML}{F2F2F2}
\definecolor{mittelgrau}{HTML}{777777}
\definecolor{dunkelgrau}{HTML}{444444}

% ═══════════════════════════════════════════════════════════════
% VARIABLEN
% ═══════════════════════════════════════════════════════════════

\newcommand{\thema}{Objetos de clase \& hay}
\newcommand{\klasse}{7}
\newcommand{\maxpunkte}{33}

% ═══════════════════════════════════════════════════════════════
% KOPF- UND FUSSZEILE
% ═══════════════════════════════════════════════════════════════

\pagestyle{fancy}
\fancyhf{}
\renewcommand{\headrulewidth}{0pt}
\renewcommand{\footrulewidth}{0pt}
\cfoot{\sanstext{\footnotesize\textcolor{mittelgrau}{\thepage\,/\,\pageref{LastPage}}}}

% ═══════════════════════════════════════════════════════════════
% BOXEN
% ═══════════════════════════════════════════════════════════════

% Grammatik-Box (elegant, Linie oben)
\newtcolorbox{grammatikbox}{
    colback=hellgrau,
    colframe=hellgrau,
    boxrule=0pt,
    arc=0pt,
    left=8pt, right=8pt, top=8pt, bottom=6pt,
    borderline north={1.5pt}{0pt}{dunkelgrau}
}

% Merk-Box
\newtcolorbox{merkbox}{
    colback=hellgrau,
    colframe=hellgrau,
    boxrule=0pt,
    arc=0pt,
    left=8pt, right=8pt, top=8pt, bottom=6pt,
    borderline north={1.5pt}{0pt}{akzent}
}

% Regel-Tag
\newcommand{\regel}[1]{%
    \tikz[baseline=(X.base)]\node[fill=akzent, text=white,
    font=\sanstext\scriptsize\bfseries, inner sep=2pt, rounded corners=1pt](X){#1};%
}

% ═══════════════════════════════════════════════════════════════
% BEFEHLE
% ═══════════════════════════════════════════════════════════════

% Spanisch: Palatino fett
\newcommand{\es}[1]{\textbf{#1}}

% Deutsch: klein, kursiv, grau
\newcommand{\de}[1]{{\small\textit{\textcolor{mittelgrau}{#1}}}}

% Lücke
\newcommand{\luecke}[1]{\rule{#1}{0.4pt}}

% Punktebox (B-Stil: klassisch)
\newcommand{\punktebox}[1]{\hfill\sanstext{\small[\_\_/#1]}}

% Aufgabentitel (konsistenter Abstand)
\newcommand{\aufgabe}[2]{%
    \needspace{4\baselineskip}%
    \vspace{12pt}%
    \noindent\sanstext{\textbf{#1}\quad#2}%
    \vspace{6pt}%
    \nopagebreak[4]%
}

% Teilaufgaben
\newlist{teil}{enumerate}{1}
\setlist[teil]{label=\alph*), leftmargin=1.8em, itemsep=3pt, parsep=0pt, topsep=3pt}

\setlength{\parindent}{0pt}
\setlength{\parskip}{4pt}
\setlength{\columnsep}{24pt}

% ═══════════════════════════════════════════════════════════════
% DOKUMENT
% ═══════════════════════════════════════════════════════════════

\begin{document}

% ═══════════════════════════════════════════════════════════════
% HEADER (D-Style, verfeinert)
% ═══════════════════════════════════════════════════════════════

\noindent
\begin{tikzpicture}[remember picture, overlay]
    \fill[kopfzeile] (current page.north west) rectangle +(\paperwidth, -1cm);
    \node[anchor=west, text=white, font=\sanstext\large\bfseries]
        at ([xshift=1.4cm, yshift=-0.5cm]current page.north west) {\thema};
    \node[anchor=east, text=white, font=\sanstext\small]
        at ([xshift=-1.4cm, yshift=-0.5cm]current page.north east)
        {Spanisch\,·\,Klasse \klasse\,·\,Arbeitsblatt};
\end{tikzpicture}

\vspace{0.5cm}

% --- Name/Datum ---
\noindent
\sanstext{\small\textbf{Name:}} \luecke{5cm} \hfill
\sanstext{\small\textbf{Datum:}} \luecke{2.5cm}

\vspace{10pt}

% ═══════════════════════════════════════════════════════════════
% EINLEITUNG
% ═══════════════════════════════════════════════════════════════

Was gibt es alles in deinem Klassenraum? In dieser Lektion lernst du, wie du auf Spanisch beschreibst, was es in einem Raum gibt. Dafür brauchst du das wichtige Wort \es{hay} -- \de{es gibt}.

\vspace{6pt}

% ═══════════════════════════════════════════════════════════════
% GRAMMATIK-ÜBERSICHT: HAY
% ═══════════════════════════════════════════════════════════════

\begin{grammatikbox}
\begin{multicols}{2}

% --- hay Grundregel ---
{\large\bfseries hay} \de{es gibt}

\vspace{4pt}

\regel{unveränderlich} -- Singular \& Plural!

\vspace{6pt}

{\small
\begin{tabular}{@{}ll@{}}
hay + un/una & \de{Hay una mesa.} \\
hay + Zahl & \de{Hay 25 sillas.} \\
hay + muchos/as & \de{Hay muchos libros.} \\
\end{tabular}}

\columnbreak

% --- Wichtige Regel ---
{\large\bfseries Achtung}

\vspace{4pt}

\regel{KEIN el/la nach hay!}

\vspace{6pt}

{\small
\begin{tabular}{@{}ll@{}}
\textcolor{akzent}{$\times$} & *Hay \textit{el} libro. \\
\textcolor{akzent}{$\times$} & *Hay \textit{la} mesa. \\
\checkmark & Hay \textbf{un} libro. \\
\end{tabular}}

\end{multicols}
\end{grammatikbox}

\vspace{4pt}

% ═══════════════════════════════════════════════════════════════
% AUFGABEN (zweispaltig)
% ═══════════════════════════════════════════════════════════════

\begin{multicols}{2}

% --- AUFGABE 1 ---
\aufgabe{1}{Ordne die Vokabeln zu.} \punktebox{8}

{\small Verbinde spanisch -- deutsch.}

\vspace{6pt}

{\small
\begin{tabular}{@{}p{3.2cm}@{\hspace{0.3cm}}p{3.2cm}@{}}
1. \es{la mesa} & a) der Stuhl \\[0.6em]
2. \es{la silla} & b) das Buch \\[0.6em]
3. \es{la pizarra} & c) der Tisch \\[0.6em]
4. \es{la ventana} & d) die Tafel \\[0.6em]
5. \es{el libro} & e) das Heft \\[0.6em]
6. \es{el cuaderno} & f) das Fenster \\[0.6em]
7. \es{el bolígrafo} & g) die Tür \\[0.6em]
8. \es{la puerta} & h) der Kuli \\
\end{tabular}}

\vspace{6pt}

{\footnotesize\textit{Lösung:} 1--\luecke{0.4cm}, 2--\luecke{0.4cm}, 3--\luecke{0.4cm}, 4--\luecke{0.4cm}, 5--\luecke{0.4cm}, 6--\luecke{0.4cm}, 7--\luecke{0.4cm}, 8--\luecke{0.4cm}}

\vspace{8pt}

% --- AUFGABE 2 ---
\aufgabe{2}{Ergänze \es{el} oder \es{la}.} \punktebox{6}

\begin{teil}
    \item \luecke{1.2cm} mesa
    \item \luecke{1.2cm} libro
    \item \luecke{1.2cm} pizarra
    \item \luecke{1.2cm} cuaderno
    \item \luecke{1.2cm} silla
    \item \luecke{1.2cm} mapa \quad {\footnotesize\textit{(Achtung!)}}
\end{teil}

\columnbreak

% --- AUFGABE 3 ---
\aufgabe{3}{Bilde Sätze mit \es{hay}.} \punktebox{6}

{\small Schau dich im Klassenraum um.}

\vspace{4pt}

{\footnotesize\textit{Beispiel:} En mi clase hay una pizarra.}

\vspace{6pt}

\begin{teil}
    \item En mi clase hay \luecke{4cm}
    \item En mi clase hay \luecke{4cm}
    \item En mi clase hay \luecke{4cm}
\end{teil}

\vspace{8pt}

% --- AUFGABE 4 ---
\aufgabe{4}{Fülle die Lücken.} \punktebox{5}

{\small Ergänze \es{hay}, \es{un}, \es{una} oder eine Zahl.}

\vspace{6pt}

En mi clase \luecke{1.5cm} muchas mesas.

\vspace{4pt}

\luecke{1.5cm} veinticinco sillas.

\vspace{4pt}

En la pared \luecke{1.2cm} \luecke{1.2cm} mapa de España.

\vspace{4pt}

En mi mochila \luecke{1.2cm} \luecke{1.2cm} libro de español y \luecke{1.2cm} cuadernos.

\end{multicols}

% ═══════════════════════════════════════════════════════════════
% AUFGABE 5: Volle Breite
% ═══════════════════════════════════════════════════════════════

\vspace{2pt}
\noindent\textcolor{hellgrau}{\rule{\textwidth}{1.5pt}}
\vspace{2pt}

\aufgabe{5}{Beantworte die Fragen.} \punktebox{4}

\begin{multicols}{2}

\textbf{a)} \es{¿Qué hay en tu mochila?}

\vspace{4pt}
\luecke{6cm}

\columnbreak

\textbf{b)} \es{¿Cuántas ventanas hay en tu clase?}

\vspace{4pt}
\luecke{6cm}

\end{multicols}

\vspace{6pt}

% ═══════════════════════════════════════════════════════════════
% AUFGABE 6: Bonus
% ═══════════════════════════════════════════════════════════════

\aufgabe{Bonus}{Beschreibe deinen Klassenraum.} \punktebox{4}

{\small Schreibe 4 Sätze. Benutze verschiedene Strukturen mit \es{hay}.}

\vspace{6pt}

\noindent\fbox{\parbox{\dimexpr\textwidth-2\fboxsep-2\fboxrule}{%
\vspace{2.5cm}
}}

% ═══════════════════════════════════════════════════════════════
% MERKKASTEN
% ═══════════════════════════════════════════════════════════════

\vspace{8pt}

\begin{merkbox}
\sanstext{\textbf{Merke: hay = es gibt}}

\vspace{4pt}

\begin{multicols}{2}

{\small
\begin{tabular}{@{}ll@{}}
\es{hay + un/una:} & Hay \textbf{una} mesa. \\[0.3em]
\es{hay + Zahl:} & Hay \textbf{cinco} sillas. \\[0.3em]
\es{hay + muchos/as:} & Hay \textbf{muchos} libros. \\
\end{tabular}}

\columnbreak

{\small\textbf{NICHT:} *Hay \textit{el} libro. \quad *Hay \textit{la} mesa.

\vspace{3pt}

(Kein bestimmter Artikel nach hay!)}

\end{multicols}
\end{merkbox}

% ═══════════════════════════════════════════════════════════════
% FOOTER
% ═══════════════════════════════════════════════════════════════

\vfill

\noindent\rule{\textwidth}{0.5pt}

\vspace{4pt}

\hfill\sanstext{\textbf{Punkte:}} \luecke{1.2cm} / \maxpunkte

\end{document}
